\section{Design Dimensions}
\label{sec:dimensions}

We propose six design dimensions for a metadata-hiding group registry:
the anchor chain, the SNARK system, the curve and in-circuit hash, the
post-quantum stance, the trusted-setup model, and the verification
host. The dimensions are not mathematically independent --- a curve
choice constrains the available SNARK family, and the verification
host axis constrains the curve --- but each imposes a binding
constraint that the others do not subsume. We treat them as a
decomposition of the bundled triple (curve, SNARK, chain) that
existing literature most often presents, and devote
\S\ref{sec:dimensions-deps} to the cross-axis dependencies that
remain.

\subsection{Anchor chain}
\label{sec:dim-chain}

The anchor chain is the public ledger that hosts the registry contract
and stores $\sigma_t$. It determines validator-set composition, the
chain's censorship-resistance posture, transaction-fee economics, and
the menu of available host functions. We consider Stellar (Soroban),
Ethereum L1, Ethereum L2 rollups (Base, Optimism, Arbitrum, zkSync,
StarkNet), Solana, Aztec, Mina, Bitcoin with inscription-style data
attachment, and a no-chain configuration that uses an availability
layer (Celestia, EigenDA) without execution. The no-chain
configuration is meaningful only if the predicate that admits state
transitions can be enforced off-chain by a quorum of operators; it
trades chain-level finality for unbounded throughput and is
out-of-scope for the deployment profiles we recommend in
\S\ref{sec:recommendations}.

\subsection{SNARK system}
\label{sec:dim-snark}

The SNARK system is the proof construction itself. We consider
Groth16~\cite{groth-2016-onpairing}, PLONK~\cite{gabizon-2019-plonk}
and its UltraPlonk variant, Halo and Halo2~\cite{bowe-2019-halo},
transparent STARKs~\cite{ben-sasson-2018-stark},
Plonky2~\cite{plonky2-2022} and Plonky3~\cite{plonky3-spec} as
hash-based instantiations targeting low-latency proving,
Nova~\cite{kothapalli-2022-nova} and its descendants as folding-scheme
recursion engines, and Spartan~\cite{setty-2020-spartan} as a
transparent-setup option without FRI. The choice fixes the proof's
asymptotic and concrete sizes, the prover's working-memory and
parallelism profile, and the post-quantum class
(\S\ref{sec:dim-pq}).

\subsection{Curve and in-circuit hash}
\label{sec:dim-curve}

For pairing-based and curve-based SNARKs, the elliptic curve and the
hash function chosen for in-circuit use jointly determine R1CS or
custom-gate constraint counts and the on-chain verifier's primitive
dependencies. We consider BLS12-381~\cite{bowe-2017-bls12381},
BN254~\cite{barreto-2005-bn254} (with the security re-evaluation
of~\cite{kim-barbulescu-2016} as a load-bearing footnote), the
BLS12-377 / BW6-761 cycle used for SNARK-of-SNARK constructions, the
Pallas/Vesta cycle (``Pasta'')~\cite{hopwood-2020-pasta} used by Halo2
and Mina, and the Goldilocks 64-bit prime field used by Plonky2/3.
The in-circuit hash is the second half of the pair: Poseidon and
Poseidon2 are common across pairing-based and FRI-based systems,
Rescue-Prime appears in STARK-friendly contexts, Pedersen hashing
predates Poseidon and is still in use in Sapling-derived constructions,
and SHA-2 family hashes appear when the constraint cost is acceptable
or unavoidable (e.g.\ for outer binding to chain-level data).

\subsection{Post-quantum stance}
\label{sec:dim-pq}

We classify configurations into three post-quantum classes:
\begin{enumerate}
  \item \textbf{Pairing-based}: soundness reduces to discrete-log on a
    pairing-friendly curve. Broken under a sufficiently large quantum
    computer running Shor's algorithm. Includes Groth16 over
    BLS12-381 / BN254 / BLS12-377, KZG-based PLONK and Halo2
    instantiations.
  \item \textbf{Classical-EC, non-pairing}: soundness reduces to
    discrete-log on an elliptic curve without pairings. Also broken
    under Shor. Includes IPA-based commitments
    (Halo2 over Pasta, Spartan-on-EC).
  \item \textbf{Hash-based}: soundness reduces to collision-resistance
    of an underlying hash. Conjectured post-quantum-secure under the
    standard model assumption that the hash remains
    collision-resistant against quantum adversaries; concretely this
    is Grover-resistance with appropriate output sizing. Includes
    STARKs and Plonky2/3.
\end{enumerate}
A fourth class --- lattice-based SNARKs --- is emerging in the
literature and is not yet represented in any deployed configuration we
survey. We flag it in \S\ref{sec:open}.

The three-class division collapses several finer distinctions; in
particular, ``conjectured PQ-secure'' is doing real work, because the
Grover speedup interacts with Fiat-Shamir security amplification in
ways that change concrete parameter choices~\cite{don-2019-fiat-shamir}.
The PQ-posture comparison table (\S\ref{sec:analysis}) records the
specific assumption each configuration relies on rather than the
class label alone.

\subsection{Trusted setup}
\label{sec:dim-setup}

Trusted-setup posture is a function of the SNARK choice but is the
operationally most-visible decision and so deserves its own axis. We
distinguish:
\begin{enumerate}
  \item \textbf{Per-circuit setup}: each relation requires its own
    Phase-2 ceremony. Groth16 is the canonical example. The ceremony
    is reusable across deployments of the same relation but not across
    relation changes; relation upgrades require a fresh ceremony.
  \item \textbf{Universal updatable setup}: a single Phase-1 ceremony
    is shared across all relations of a given size; new contributors
    can extend the ceremony without restarting. PLONK and Marlin are
    canonical examples.
  \item \textbf{Transparent setup}: no ceremony. STARKs, Plonky2/3,
    Halo2-with-IPA, and Spartan are transparent. The trade-off is
    typically larger proof size and (for FRI-based systems) higher
    verifier work.
\end{enumerate}
The number of contributors required for a given setup to be
soundness-bearing is one; trustworthiness against the trapdoor is
maximised when at least one contributor in fact destroyed their
toxic waste. Multi-contributor MPC ceremonies make this more credible
but do not change the formal property.

\subsection{Verification host}
\label{sec:dim-host}

The verification host axis is whether the chain provides native
support for the SNARK's verifier primitives or whether the verifier
runs in the chain's general-purpose bytecode. Three regimes:

\begin{enumerate}
  \item \textbf{Native precompile / host function}: the chain exposes
    a built-in callable that performs (e.g.) a BLS12-381 pairing or
    a BN254 G1 multi-scalar multiplication at gas costs orders of
    magnitude lower than a bytecode implementation. Examples:
    Ethereum's BN254 precompile (EIP-197), Stellar Protocol 22's
    BLS12-381 host functions~\cite{stellar-cap-0058},
    Mina's built-in Pickles verifier.
  \item \textbf{Pure bytecode}: the verifier is a contract
    implementation in the chain's VM language with no native
    acceleration for SNARK primitives. Verification cost scales with
    the bytecode interpreter's per-operation overhead.
  \item \textbf{Verifier-as-rollup}: the chain itself is a SNARK-rollup
    whose execution model is the SNARK verifier, so application-level
    proofs are checked by the same machinery that validates the
    chain. Aztec is the canonical example.
\end{enumerate}

This axis is the binding constraint that makes the
``Stellar~+~Plonky3~+~FRI'' configuration in \S\ref{sec:configurations}
hypothetical. Stellar Protocol 22 added BLS12-381 host functions but
no FRI host functions; a transparent-setup PQ-secure verifier on
Stellar would run in pure bytecode and incur orders-of-magnitude
higher gas. The same constraint forces Ethereum L1 deployments toward
BN254 over BLS12-381 in the absence of EIP-2537.

\subsection{Cross-axis dependencies}
\label{sec:dimensions-deps}

The six axes carry the following hard dependencies:

\begin{itemize}
  \item \emph{SNARK $\to$ trusted-setup}. Groth16 forces per-circuit;
    PLONK forces universal-updatable; STARK / Plonky2 / Plonky3 force
    transparent. The trusted-setup axis is reducible in theory, but
    operationally distinct because the ceremony's
    multi-contributor-MPC posture is independent of the SNARK choice.
  \item \emph{SNARK $\to$ post-quantum stance}. Pairing-based SNARKs
    are not PQ; hash-based SNARKs are PQ-conjectured. PQ posture is
    therefore reducible to (SNARK, curve), but again the
    operationally-relevant question --- ``which Grover security level
    does this configuration target?'' --- is one further refinement.
  \item \emph{Anchor chain $\to$ verification host}. The available
    host functions are a property of the chain. Stellar gives
    BLS12-381 pairings; Ethereum gives BN254; Mina gives Pickles;
    Aztec gives UltraPlonk-as-execution-model.
  \item \emph{Verification host $\to$ curve / hash}. The available
    host functions transitively constrain the curve choice when
    in-bytecode verification is uneconomical.
\end{itemize}

The remaining cells of the (axis $\times$ axis) dependency matrix are
soft. The taxonomy's value is that it makes the dependency direction
explicit: a designer who fixes the chain for non-cryptographic reasons
(team familiarity, regulatory posture, fee economics) is implicitly
fixing four of the six axes via the dependency chain
chain~$\to$~host~$\to$~curve~$\to$~SNARK~$\to$~setup, and the only
remaining freedom is the in-circuit hash choice.

Figure~\ref{fig:axes} renders the six axes with example configurations
from \S\ref{sec:configurations} placed at their corresponding points;
the dependency arrows above are visualised as directed edges between
axes.

% Figure 1: Design-space axes diagram for §3.
% Six axes radiating from a central point; representative configurations
% from §5 marked as labeled points.
%
% TODO: design after §3 prose is drafted. Stub renders nothing.
\begin{figure}[t]
  \centering
  % \begin{tikzpicture}
  %   ...
  % \end{tikzpicture}
  \fbox{\parbox{0.9\columnwidth}{\centering\vspace{4em}TODO Figure 1:
    six-axis taxonomy diagram\vspace{4em}}}
  \caption{Six orthogonal design axes for metadata-hiding group
    registries with SNARK-gated state changes. TODO redraw.}
  \label{fig:axes}
\end{figure}

